\section{Direct C222 equations with one oriented Pell cap}
\label{sec:direct-cap}
\label{sec:arithmetic}

The direct proof-calculus route now gives two natural quartics and two
integer quartics.  Their coefficients are small enough to be supplied as
complete collected lists.  The integer constructions use a new bounded
congruence that recovers the signs and indices of fifteen Pell pairs at
once.  This section proves the new arithmetic and its composition with
the inherited numerical C222 calculus.  It does not assert a completed
Lean or Rocq derivation.

\begin{center}
\begin{tabular}{lrrrr}
\hline
Domain and variables & Rows & Degree & Terms & Height\\
\hline
$\mathbb N^{240}$ &188&4&25,987&90,589,484\\
$\mathbb N^{239}$ &187&4&360,097&491,869,488\\
$\mathbb Z^{381}$ &225&4&13,179&2,445,790,974,704\\
$\mathbb Z^{391}$ &235&4&10,764&80,496,640\\
\hline
\end{tabular}
\end{center}

Here $\mathbb N=\{0,1,2,\ldots\}$. Each equation is the sum of the
squares of every quadratic residual in its corresponding literal JSON
file. Every declared variable occurs, and the collected single
polynomial has the exact degree and coefficient height shown.  The
240-variable form improves the inspected 261-variable direct natural
form by 21 variables, six rows and 5,692 terms, at a 5.9\% increase in
height. The 381-variable integer form saves 61 variables and 65 rows
against the inspected 442-variable, 290-row direct integer form, but
uses more terms and larger coefficients. These are distinct tradeoffs;
the comparison is confined to the direct C222 architecture.

\subsection{The inherited source and the new composition}

The numerical proof calculus C222 has an inherited, machine-independent
refutation calibration through the named endpoints
\nolinkurl{c222_iff_zf_refutation} and
\nolinkurl{zf_refutation_iff_zfc_refutation}.  The unchanged source
uses beta coding, a nonnegative Horn-flow polynomial, amplified
prime-power CRT, and exact binomial digits.  These complete ordinary
proofs and literal source files are retained in the deliverable.
The new construction does not simulate a K215 execution or instantiate
Jones's universal polynomial.

The incoming shared-binomial construction~\cite{binomial} already uses one binomial
for both a factorial fixed point and a reciprocal CRT product, and
one common Pell residue for eleven binomial numerators.  It retains
six exact power requests but fifteen complete indexed-Pell certificates.
The separate incoming shared-cap construction~\cite{cap} replaces many indexed
certificates by one cap, assuming every requested index belongs to its
power-size frame.  Their combination requires a changed frame: the
modular numerator indices must be included even though their enormous
power outputs are absent.  The compatibility of that enlarged frame
with the cofinal common-root parameter is proved below.

The ordinary source has $t=t_0+16$, a proposed natural factorial $F$,
beta parameters $x,y$, and nine positive encoded CRT representatives.
Set
\[
 B=b_0+2,\quad R=2+t+x+y+F,\quad K=2R,
 \qquad V_0=(F^2+1)^t,\quad N=B^K,\quad q=N^K.
\]
The six exact powers consist of $V_0,N,q$ and
\[
 D_t=q^t,\qquad D_y=q^{y+1},\qquad D_x=q^{x+1}.
\]
The adjacent denominator is the ordinary product
$D_{t+1}=qD_t$.  The one exact binomial digit will identify
$B=\binom{F^2+t}{t}$.  There are ten additional range divisibilities:
the nine independent CRT ranges and the derived-label range.
Every upper index is positive before any power or factorial is
interpreted.  The two raw ceiling rows are
\begin{equation}
 V_0+B=BF+r+1,\qquad B=r+s+1,
 \qquad r,s\in\mathbb N.
 \label{eq:direct-ceiling}
\end{equation}

Let $\mathcal A$ be the set of distinct atoms occurring as bases,
exponents or outputs of the six exact requests, together with the
upper indices of the eleven modular requests.  Its literal size is 23.
The natural construction defines
\begin{equation}
 W=2+\sum_{a\in\mathcal A}a.
 \label{eq:direct-natural-frame}
\end{equation}
Its fifteen distinct requested indices $k_i$ are positive atoms in
$\mathcal A$. Thus $k_i<W$ and $m=15\le W-2$.  No private Pell
coordinate occurs in this definition.

\subsection{A cofinal base above the enlarged frame}

For $a\ge2$, write $(X_j(a),Y_j(a))$ for the Pell sequence defined by
\[
 X_0=1,\ Y_0=0,\qquad
 X_{j+1}=aX_j+(a^2-1)Y_j,\quad Y_{j+1}=X_j+aY_j.
\]
Its elementary properties used below are Pell classification,
strict increase of $Y_j$, and
\[
 X_j(a)\equiv1\pmod{a-1},\qquad
 Y_j(a)\equiv j\pmod{a-1}.
\]
Retain natural witnesses $L,\sigma,Z$ with $Z\ge1$ and
\begin{equation}
 L^2-\bigl((W+1)^2-1\bigr)(WZ)^2=1,
 \qquad A=L+\sigma.
 \label{eq:direct-large-base}
\end{equation}
Pell classification writes $L=X_e(W+1)$ and $WZ=Y_e(W+1)$.
Since $WZ>0$, $e\ge1$; reduction modulo $W$ gives $W\mid e$.
Consequently
\[
 A\ge L\ge X_W(W+1)\ge(W+1)^W.
\]
Conversely $e=W$ supplies $L,Z$, while $A$ remains free to be any
larger integer. This freedom is essential.

Put $b=q+1$ and $M=qD_{t+1}D_xD_y$. The common-root condition is
\begin{equation}
 2Ab=b^2+1+Mh_0.
 \label{eq:direct-common-root}
\end{equation}
For the intended semantic powers, $M$ is a positive power of $q$.
Thus $\gcd(b,M)=1$ and $\gcd(2b,M)\in\{1,2\}$. In the second case
$b$ is odd and $b^2+1$ is even. The linear congruence for $A$ is
soluble, and its solution class has arbitrarily large positive members.
Choose one above $L$ and large enough that $h_0\ge0$. Hence the
enlarged frame and the common root admit simultaneous witnesses.

\subsection{The natural shared cap}

One complete inherited indexed-Pell certificate at index $W$ gives
$V_c=Y_W(A)$. For clarity its polynomial interface is
\begin{align*}
 X_c^2-(A^2-1)V_c^2&=1,& U^2-(A^2-1)V^2&=1,\\
 S^2-(B_c^2-1)T^2&=1,& B_c&=1+4V_c(d+1),\\
 V&=V_c^2(e+1),&B_c&=A+Uf,\\
 S&=X_c+Ug,&T&=W+4V_ch.
\end{align*}
All cap witnesses here are natural. Its usual condition $W\le V_c$
is supplied by the aggregate row below. The indexed-Pell criterion and
its constructive completeness are retained dependencies.

For each of the fifteen other pairs retain only
\[
 x_i^2-(A^2-1)y_i^2=1,
\]
and impose
\begin{equation}
 V_c=W+1+s_c+\sum_i y_i,
 \qquad \sum_i(y_i-k_i)^2=(A-1)Q,
 \quad s_c,Q\in\mathbb N.
 \label{eq:direct-natural-cap}
\end{equation}
Classification gives $y_i=Y_{j_i}(A)$ for $j_i\ge0$.
The first row implies $y_i<V_c$, hence $j_i<W$.
The second implies that $A-1$ divides
$E=\sum_i(j_i-k_i)^2$. But
\[
 0\le E\le m(W-1)^2<W^3<A-1,
\]
so every $j_i=k_i$. This bounds the decoded indices, not the enormous
ordinates.

For completeness choose the canonical pairs. Each $y_i-k_i$ is a
nonnegative multiple of $A-1$. The quotient $Q$ is therefore natural.
Moreover $\sum_i y_i\le mY_{W-1}(A)$ and
$Y_W(A)>(2A-1)Y_{W-1}(A)$. Since $m\le W-2$ and $A>W+1$, the
required slack $s_c=Y_W(A)-W-1-\sum_i y_i$ is nonnegative.

\subsection{The oriented cap over unrestricted integers}

The following is the new sign-removal mechanism. Let $W\ge4$,
$A-1>4W^3$, $1\le k_i<W$, $m\le W-2$, and suppose
$V_c=Y_W(A)$ is already certified. Permit $x_i,y_i,Q$ to be arbitrary
integers, retain their principal norms, and use only
\begin{align}
 V_c&=W+1+s_c+\sum_i y_i^2,&s_c&\in\mathbb N,
 \label{eq:direct-square-cap}\\
 \sum_i\bigl((y_i-k_i)^2+(x_i-1)^2\bigr)&=(A-1)Q.
 \label{eq:direct-oriented-error}
\end{align}
Then
\begin{equation}
 (x_i,y_i)=(X_{k_i}(A),Y_{k_i}(A))
 \quad\text{for every }i.
 \label{eq:direct-oriented-conclusion}
\end{equation}

Indeed signed Pell classification gives
$x_i=\varepsilon_iX_{j_i}(A)$ and
$y_i=\delta_iY_{j_i}(A)$ with independent signs in $\{1,-1\}$.
Equation~\eqref{eq:direct-square-cap} implies $y_i^2<V_c$ and hence
$|y_i|<V_c$, using $|u|\le u^2$ for nonzero integers.
Thus $0\le j_i<W$. Reducing~\eqref{eq:direct-oriented-error}
modulo $A-1$ shows that this modulus divides
\[
 E=\sum_i\bigl((\delta_i j_i-k_i)^2+(\varepsilon_i-1)^2\bigr).
\]
Every summand is nonnegative, and
\[
 0\le E\le m\bigl(4(W-1)^2+4\bigr)<4W^3<A-1.
\]
Therefore $E=0$. It follows that $\varepsilon_i=1$ and
$\delta_i j_i=k_i>0$, proving~\eqref{eq:direct-oriented-conclusion}.
The extra $(x_i-1)^2$ terms are essential for recovering abscissa signs.
There is no sign assumption on the quotient $Q$.

For constructive completeness assume also $2k_i\le W-1$ and $A>W+1$.
Choose all pairs canonical. Their differences $x_i-1$ and $y_i-k_i$
are divisible by $A-1$, so~\eqref{eq:direct-oriented-error} supplies
a natural $Q$. The addition law gives
\[
 Y_{k_i}(A)^2<Y_{2k_i}(A)\le Y_{W-1}(A).
\]
The same recurrence estimate as before makes
$Y_W(A)-\sum_iY_{k_i}(A)^2>W+1$, giving the required natural slack.

The integer source replaces~\eqref{eq:direct-natural-frame} by
\begin{equation}
 W=2+\sum_{a\in\mathcal A}a^2.
 \label{eq:direct-square-frame}
\end{equation}
Even with signed power outputs this bounds every semantic atom in
absolute value by $W$. All fifteen requested indices are positive
before power interpretation, so $m\le W-2$ and
$W\ge k_i^2+2\ge2k_i+1$. The literal $t\ge16$ gives $W\ge258$.
Equation~\eqref{eq:direct-large-base} still yields
\[
 A-1\ge(W+1)^W-1\ge(W+1)^4-1>4W^3.
\]
Thus the square frame proves every hypothesis of the oriented cap.
If zero indices were permitted in a general formulation, a separate
count bound would be needed; they do not occur here.

\subsection{Sequential recovery of power and quotient signs}

After the cap has identified every principal pair, an exact power row is
\begin{equation}
 X_k(A)+(n-A)Y_k(A)=m+T_nu,
 \qquad T_n=2An-n^2-1.
 \label{eq:direct-signed-power}
\end{equation}
Allow $m,u,T_n$ initially to be signed, but suppose $n>0$ is already
known. The square frame gives $1\le n,k<W$ and $|m|<W$.
The polynomial defining $T_n$ is increasing on $1\le n<W$, so
\[
 T_n\ge2A-2>W^W+W,
 \qquad 0<n^k\le(W-1)^{W-1}<W^W.
\]
The usual Pell congruence and~\eqref{eq:direct-signed-power} imply
$T_n\mid m-n^k$. Since $|m-n^k|<T_n$, they force $m=n^k>0$.
The quotient is then nonnegative because the Pell expression is positive
and $m$ is its least positive residue.

The six calls are ordered with bases $F^2+1$, $B$, $N$, and then $q$
three times. The first two are independently positive; their established
outputs license the later bases. The sign argument is consequently
noncircular. It removes the four modulus sign guards and the two output
sign guards that actually remained in the earlier finite census.

The common-root, CRT-index and local-evaluation quotients are needed
only as integer divisibility witnesses for soundness. Their natural signs
follow later from the positive denominators and nonnegative numerators.
The modular-digit quotients are also unrestricted premises: after
recovery the Pell residue is positive and the selected low residue is
least nonnegative, so those quotients are automatically nonnegative.

\subsection{Carry bounds and source calibration}

Equation~\eqref{eq:direct-common-root} gives, for every $n\ge0$,
\[
 X_n(A)+(q+1-A)Y_n(A)\equiv(q+1)^n\pmod M.
\]
This follows by substituting $q+1-A$ for the square root in the Pell
polynomial identity. The resulting expression is positive. Every
required exact-digit modulus $qd$, and every range modulus $Nd$, divides
$M$; the adjacent denominator $D_{t+1}=qD_t$ is included explicitly.
The unchanged remainder rows therefore give the intended digits once
their carry bounds have been established.

The bootstrap is independent of factorial meaning. From $B\ge2$ and
the recovered exact powers, $q\ge2^{4R^2}$. For $0\le j\le t$,
\[
 \binom{F^2+t}{j}\le(F^2+t)^t
 \le(R^2+R)^t\le2^{2Rt}<2^{2R^2}<q.
\]
Thus the exact digit identifies $B=\binom{F^2+t}{t}$.
Equation~\eqref{eq:direct-ceiling} then forces $F>0$ and its inherited
square-radius fixed-point theorem gives $F=t!$.
The retained source proof subsequently establishes every other carry
bound, the reciprocal moduli
$d_i=1+F^2/i$ and their pairwise coprimality, the amplified CRT ranges,
and the finite Horn refutation.  None of these meanings is used before
the preceding power and binomial recovery.

Conversely a finite C222 refutation supplies a padded horizon, its true
factorial, beta parameters, CRT representatives and semantic powers.
Those data determine either enlarged frame before any private Pell
witness. Choose a large compatible $A$ by
\eqref{eq:direct-large-base}--\eqref{eq:direct-common-root}, construct
the one indexed cap and canonical principal pairs, and use the explicit
slacks above. All power and digit quotients then exist. This proves
both directions of the arithmetic source equivalence.

\subsection{Guard graphs, finite identities, and the endpoint}

The integer source retains 52 explicit nonnegativity predicates.
For each selected coordinate $z$, replace its sign condition by the
quadratic row
\[
 z=a^2+b^2+c(c+1),\qquad a,b,c\in\mathbb Z.
\]
This expression is nonnegative and represents every natural number:
apply Legendre's three-square theorem to $4z+1$ and separate the unique
odd square. Every further required sign follows from 174 literal row
consequences, each checked using already known signs or an exact affine
square. The fifteen index-positivity proofs use only the pre-power
source witnesses. The proof never silently treats an exempt signed
coordinate as natural.

The 349-coordinate raw quadratic source with its three changed rows and
52 guard triples has 505 integer variables and 349 residuals. Integral
monic substitution removes 124 coordinates and rows, leaving the stated
381-variable, 225-row quartic. Restricting affine support instead gives
the 391-variable lower-height form. No unconstrained coordinate is
deleted. All inverse substitutions and complete expansions are supplied.

The independent checker imports no source producer, sign-inference
routine or optimizer. It reconstructs the three changed rows, compares
all unchanged rows, checks every guard and sign step, proves the exact
pre-power positivity cut, replays every monic restoration, and compares
every collected quartic coefficient. Separate altered-coefficient and
altered-certificate controls are rejected. These finite checks establish
the emitted arithmetic syntax; the unbounded proof and inherited logic
remain the ordinary arguments just identified.

Negating the resulting equivalence between a polynomial zero and a
finite ZFC refutation gives, for each of the four supplied polynomials,
\[
 \mathrm{ZFC}\ \vdash\
 \bigl(\text{no zero in the stated domain}\bigr)
 \ \longleftrightarrow\ \operatorname{Con}(\mathrm{ZFC}).
\]
